% !TeX root = ../tesis.tex

%-----------------------------------------------------------------------
\section{Objetivos de la Investigación}
\label{sec:objetivos}
%-----------------------------------------------------------------------

\subsection{Objetivo General}
\label{subsec:obj-general}

Diseñar y documentar una metodología replicable basada en herramientas de
código abierto para el diagnóstico cuantitativo de redes de transporte
multimodal urbano ---sustentada en el procesamiento de datos \gtfs{} y el
modelado de grafos matemáticos--- con el fin de caracterizar las
vulnerabilidades funcionales, la distribución de cobertura y las desigualdades
territoriales de acceso en la red de la \zmvm{}; y demostrar su capacidad
analítica mediante la simulación del impacto hipotético de la implementación
de corredores anillares de alta capacidad en los cuatro cuadrantes del Anillo
Periférico de la \cdmx{}.

\subsection{Objetivos específicos}
\label{subsec:obj-especificos}

\begin{enumerate}
\item Sistematizar y procesar los feeds de datos abiertos (\gtfs{}) de la
    totalidad de la red de Movilidad Integrada de la \cdmx{} y de los sistemas
    de transporte del Estado de México (Mexibús, Mexicable) mediante una
    herramienta de adquisición y normalización de código abierto, con el fin
    de construir un grafo multimodal georreferenciado como insumo base del
    análisis.

\item Calcular el conjunto de indicadores de eficiencia y cobertura de la red
    de transporte multimodal de la \zmvm{} ---incluyendo el Índice de Ruta
    Directa ($DI$), el Tiempo de Viaje Promedio ($T$), el Nivel de Cobertura
    ($C$) y la Fuerza Capilar ($FC$)--- mediante la cadena de herramientas
    de código abierto desarrollada, con el fin de obtener un diagnóstico
    cuantitativo de la configuración topológica actual de la red.

\item Identificar los nodos críticos y vulnerabilidades estructurales de la
    red de la \zmvm{} mediante el cálculo de la Centralidad de Intermediación
    ($B(v)$) y la Robustez Geométrica ($\Delta E$), con el fin de localizar
    los puntos neurálgicos cuya falla comprometería la conectividad global
    del sistema y fundamentar la pertinencia de alternativas perimetrales.

\item Identificar las zonas periféricas con mayor déficit de cobertura
    estructurada y conectividad lateral en la \zmvm{} a partir de los
    resultados de Nivel de Cobertura ($C$) y Fuerza Capilar ($FC$),
    con el fin de delimitar las áreas de intervención prioritaria para
    la evaluación del escenario hipotético anillar.

\item Simular el impacto hipotético de la incorporación de corredores anillares
    de alta capacidad en los cuatro cuadrantes del Anillo Periférico mediante
    la comparación de los indicadores $DI$, $T$ y $B(v)$ entre la red actual
    y la red ampliada, con el fin de demostrar la capacidad analítica de la
    metodología para evaluar escenarios alternativos de planeación.

\end{enumerate}

\subsection{Pregunta de Investigación}
\label{subsec:pregunta}

¿En qué medida una metodología replicable basada en herramientas de código
abierto y modelado de grafos matemáticos permite caracterizar las
vulnerabilidades funcionales y brechas de cobertura de la red de transporte
multimodal de la \zmvm{}, y demostrar analíticamente el impacto hipotético
de corredores anillares en los cuatro cuadrantes del Anillo Periférico sobre
la eficiencia global del sistema?


De manera complementaria, se plantean las siguientes preguntas secundarias:

\begin{itemize}
    \item ¿De qué manera la alta dependencia de nodos centrales (\textit{hubs})
    en la Ciudad Central y el Primer Contorno compromete la robustez del
    sistema ante perturbaciones operativas \citep{semovi2020diagnostico}?

    \item ¿Cómo contribuye la brecha de conectividad lateral en el Tercer
    Contorno al incremento de los tiempos de traslado para los usuarios
    provenientes del Cuarto Contorno \citep{semovi2020diagnostico}?

\end{itemize}

